\chapter{Проблемно-классификационный анализ задачи оценки пространственной доступности городских сервисов}
\label{ch:problem-analysis}

\section{Городские сервисы как объект анализа доступности}
\label{sec:urban-services}

\subsection{Типы городских сервисов и показатели обеспеченности}
\label{subsec:service-types}

\subsection{Спрос, предложение и пространственное распределение сервисов}
\label{subsec:demand-supply-space}

\subsection{Ограничения существующих подходов к агрегированной оценке обеспеченности}
\label{subsec:provision-limitations}

\section{Транспортная и пешеходная доступность в задачах оценки городских сервисов}
\label{sec:transport-walk-access}

\subsection{Пешеходная доступность между местом проживания, остановками и сервисами}
\label{subsec:first-last-mile}

\subsection{Доступность на основе общественного транспорта}
\label{subsec:pt-accessibility}

\subsection{Комбинированные сценарии доступности с учетом нескольких модальностей}
\label{subsec:multimodal-access}

\section{Уличная морфология как фактор неравномерности доступа}
\label{sec:street-pattern-factor}

\subsection{Классификации street pattern и топологические характеристики уличных сетей}
\label{subsec:street-pattern-taxonomy}

\subsection{Связь street pattern с first-mile и last-mile провалами доступности}
\label{subsec:street-pattern-walk-failures}

\subsection{Роль street pattern на маршрутах общественного транспорта}
\label{subsec:street-pattern-pt-paths}

\section{Методы моделирования и оптимизации размещения городских сервисов}
\label{sec:service-placement-review}

\subsection{Классические модели размещения объектов обслуживания}
\label{subsec:classical-flp}

\subsection{Модели покрытия спроса и учета неудовлетворенного спроса}
\label{subsec:coverage-demand}

\subsection{Совместная интерпретация транспортной связности и размещения сервисов}
\label{subsec:transport-service-joint}

\section{Постановка цели и задач исследования}
\label{sec:research-goal-tasks}

\subsection{Исследовательский разрыв}
\label{subsec:research-gap}

\subsection{Цель и основные задачи диссертационного исследования}
\label{subsec:goal-tasks}

\subsection{Ожидаемые научные и практические результаты}
\label{subsec:expected-results}

\section{Выводы по главе}
\label{sec:chapter1-conclusions}
